\section{普通の岩石から濃集を説明する}\label{theory:ore:origins}
鉱床は特定の元素・鉱物が周囲より濃集した地質体として考える。濃集の存在、品位、鉱量、採掘可能量、採算は別の評価段階である。品位が低くても大規模な鉱体は資源となり得る。
まず元素の供給源を考え、運搬媒体、通路、濃集条件、保存を順に追う。正マグマ性は結晶分化・集積や硫化物液の分離、熱水性は流体中の成分の沈殿、堆積性は化学沈殿や粒子運搬、風化残留は移動しにくい成分の残留を中心に整理する。石油は根源岩での生成・移動・貯留とトラップを分ける。
層状苦鉄質貫入体のクロムと白金族元素は代表的なマグマ性資源である\cite{ore50_magmatic_ref}。金銀の熱水成因が重要であることから、全マグマ性鉱床に金が絶対にないとは推論しない。自然金を含むPGE鉱床モデルもある\cite{ore50_gold_exception_ref}。
水中の化学種・pH・酸化還元状態で移動性は変わる。元素名の固定した溶けやすさ順位だけで全成因を説明しない。

\section{熱水の通路と沈殿条件}\label{theory:ore:hydrothermal}
高温の流体は金属等を運び、温度低下、圧力変化、流体混合、酸化還元やpHの変化で鉱物を沈殿させ得る。岩石の亀裂・断層は通路であり、通路を充填すれば鉱脈となる。流体の全てがマグマだけに由来するわけではない。
斑岩銅は浅所の貫入岩とその周辺に鉱染・細脈網を伴う大規模低品位の熱水系として考える。収束縁の大陸弧と海洋性島弧の両方のモデルがある\cite{ore50_porphyry_ref}。大陸弧に多いという傾向を島弧に絶対存在しないという全称に変えない。
炭酸塩岩との化学交代でCaケイ酸塩鉱物を作るスカルンは、熱だけによる再結晶と違い、成分の移動を伴う。高温から後期低温へ鉱物組合せが変化する例がある\cite{ore50_skarn_ref}。温度だけで金属種を全て決めず、母岩・流体組成・酸化還元条件も考える。

\section{黒鉱・グリーンタフ・別子型}\label{theory:ore:kuroko}
黒鉱は代表的に珪長質海底火山活動と関係する塊状硫化物鉱床で、閃亜鉛鉱、方鉛鉱、黄鉄鉱、黄銅鉱、四面銅鉱、重晶石等を含む\cite{ore50_kuroko_ref}。海底熱水からの化学的沈殿と、その下の流体通路を考える。
日本の代表的な産地は東北の中新世グリーンタフ地域にある。グリーンタフは変質した火山岩・火山砕屑岩を含む地域的呼称であり、緑色片岩という広域変成岩の岩種名とは違う。名称の「緑」だけから同じ母岩としない。
層状含銅硫化鉄鉱の別子型は、海底熱水性の硫化物が変形・変成を経たモデルで、三波川帯の例にも変成相の記録がある\cite{ore50_besshi_ref}。黒鉱の代表産状と、変成した苦鉄質・堆積岩系に産する硫化物を区別する。
年代区分の西黒沢期に「全て限られる」という原典の全称は当時の分類と例外を確認する。代表的な集中地域の証拠を、全世界の全型の排他的条件へ変換しない。

\section{プレート境界と鉱床の対応}\label{theory:ore:tectonic}
中央海嶺は発散、海溝での沈み込みは収束である\cite{usgs_plate_boundaries}。島弧・背弧でも海底火山と熱水系が発達し、黒鉱型の潜在鉱床が報告される\cite{ore50_backarc_ref}。斑岩銅も海洋性島弧・大陸弧等に関係するため、一つの境界呼称から一つの鉱床だけを決める対応表は不十分である。
平成24年度問25の「マリアナ型」に対する独立解答と公式番号の相違は問題別に保存した。母岩、伸張・圧縮、マグマ、熱水の観測を比べる必要があり、公式4を得るために黒鉱型の島弧・背弧での存在を否定しない。砂岩型ウランは堆積盆地内の透水性と酸化還元境界を考え、収束境界だけの資源としない。

\section{堆積資源と風化残留}\label{theory:ore:sedimentary}
縞状鉄鉱は珪質層と鉄に富む層の互層であり、代表モデルでもチャートと鉄鉱物層の組合せを扱う\cite{ore50_bif_ref}。鉄の酸化状態と海水での沈殿、堆積後の変成・変形を区別する。原典で銅へ置換した成因をそのまま覚えない。
漂砂では鉱物粒子を物理的に運搬・分級する。砂金の濃集と、溶存元素が還元帯で化学沈殿する砂岩型ウランを区別する。
ボーキサイトでは強い風化で他成分が溶脱し、Alを含む難溶物が相対的に残る\cite{ore50_bauxite_ref}。Al・FeとNa・K・Ca・Mgの残留・溶脱を逆転させた候補が出題される。濃度上昇は対象元素の新規供給だけでなく、分母となる他成分の減少でも起こる。
オイルシェールはケロジェンを含む細粒堆積岩である\cite{ore50_oilshale_ref}。粗粒・細粒の交換に注意する。メタンハイドレートは水の構造へガスが取り込まれる水和物で、温度圧力による安定条件をもつ。見た目のシャーベット状という説明だけで化学構造を定義しない。

\section{石炭化を年代だけで決めない}\label{theory:ore:coal}
植物由来の有機物が埋没・加熱を経て炭素に富み、水分・揮発分等が減る石炭化を考える。基本のランク順は褐炭、亜瀝青炭、瀝青炭、無煙炭である\cite{ore50_coal_ref}。石墨に近い状態までを扱う場合も、低ランクの褐炭を瀝青炭の後に置かない。
石炭化に伴うH/C・O/Cの変化は一般に減少方向となる。原子比、質量割合、揮発分、含水量は別の指標であり、分母を確認する。
日本には古第三紀の炭田が多く、主にジュラ系とする説明は適切でない\cite{ore50_japan_coal_ref}。石炭紀の代表的な湿地植生を針葉樹だけへ交換しない。
形成温度の概説に150--200℃の代表範囲があるが、ランクは最高到達温度と埋没熱史に依存する\cite{ore50_coal_temp_ref}。古生代という年代だけで全てが200℃以下と保証できない。年代と熟成度の二つを独立に扱う。

\section{地熱の熱・水・透水性}\label{theory:ore:geothermal}
地熱は地殻の放射性壊変、深部からの熱、火山近くの局所熱等を含む。熱いだけで流体を取り出せるとは限らず、流体と有効な透水経路を持つ貯留層が必要となる。亀裂・断層や封じるキャップロックが重要である\cite{ore50_reservoir_ref}。
発電では蒸気をタービンへ使う方式と、熱交換で低沸点二次媒体を蒸発させるバイナリー方式を分ける。中・低温の温水は暖房等へ直接利用でき、非火山地域にも温水系がある\cite{ore50_geothermal_ref}。
概略の取出し熱量率は \(\dot Q=\dot m c_p(T_{\rm in}-T_{\rm out})\)。流量・温度差・持続する涵養が関係し、熱の存在量と実際に取り出せる出力を区別する。バイナリーでも熱力学的効率は温度差に制約される。
地化学温度計は溶解平衡等の温度依存を利用する。平衡到達、上昇時の冷却、異なる水との混合、再平衡の有無を確認する。地温勾配の代表20--30、局所50--100℃/km以上等を全地点の固定値にしない。

\section{深海団塊の極めて遅い成長}\label{theory:ore:marine}
マンガン団塊は鉄・マンガン酸化物を主体に核の周囲へ成長し、Cu・Ni・Co等を伴う。海底の堆積が遅い場所に露出する団塊と、海山の岩盤を覆うクラストを区別する。
JOGMECの代表説明では成長速度は百万年で1mm程度である\cite{ore50_nodule_ref}。数mm／1万年への単位期間の置換は二桁以上の差になる。起源・続成作用・環境による幅を認めて代表率として読む。原典の副成分最大1--2\%という表現は全団塊の絶対上限として未確定であり、問題別に留保する。
成長の遅さ、資源量、採掘試験、商業生産、採算、環境影響の確認は別の事項である。海に鉱物があるという観測だけで商業生産中と結論しない。

\section{ウランの母岩・流体・年代・品位}\label{theory:ore:uranium}
代表的な分類に先カンブリア時代の石英礫岩、原生代の不整合関連、顕生代の砂岩型がある\cite{ore50_uranium_ref}。古い礫岩では酸素の少ない条件下で鉱物粒子の保存と漂砂濃集を考える。砂岩型では流体が溶存Uを運び、還元帯等で沈殿させる。どちらも堆積岩中だから同じ機構とはしない。
不整合関連では古い基盤と盆地堆積物の境界付近の流体鉱化、鉱脈型では地質体を切る亀裂充填、貫入岩型では岩体に伴う分散した鉱化等を考える。
IAEAの資源説明にも低・非常に低品位で大規模な貫入岩関連等の資源がある\cite{ore50_intrusive_ref}。貫入岩型の全てを高品位としない。品位が高いことと総量が大きいことは別である。
原典の4億年・23億年等の年代区切りは分類史・対象型の範囲と照合が残る。厳密な世界全例の境界として教えず、代表的地質時代と酸化環境を結び付けて理解する。

\section{学習順序}
まず濃集の機構を五種類に分け、次に母岩・鉱物・形態を対応させる。その後プレート環境と時代、温度・酸化還元条件を加える。最後に品位・鉱量・生産状況という資源評価を読む。似た名前や緑色という共通点だけで鉱床を統合せず、原典の存在命題と全称命題を区別する。
